% =============================================================================
% preamble.tex — packages, fonts, BiDi, headers/footers, bibliography.
% Engine: LuaLaTeX (assignment §13.2). Compile with: lualatex + biber + lualatex.
% =============================================================================

% --- Page geometry -----------------------------------------------------------
\usepackage[a4paper,margin=2.5cm,headheight=15pt]{geometry}

% --- Fonts + Hebrew/English bidirectional typesetting (babel, bidi=basic) -----
\usepackage{fontspec}
% babel's Unicode bidi engine keeps embedded LTR (English) runs upright inside
% RTL (Hebrew) text and reverses RTL correctly. This is the modern, robust route
% for Hebrew under LuaLaTeX; it replaces polyglossia+luabidi, whose LuaTeX bidi
% mirror-reversed every embedded English run.
\usepackage[bidi=basic]{babel}
\babelprovide[import,main]{hebrew}
\babelprovide[import]{english}

% Hebrew fonts ship with the `culmus` package. This is the ONLY block to change
% for fonts — e.g. swap "David CLM" for "Frank Ruehl CLM" or "Noto Sans Hebrew".
\babelfont{rm}[Script=Hebrew]{David CLM}       % Hebrew serif (main body)
\babelfont{sf}[Script=Hebrew]{Nachlieli CLM}   % Hebrew sans
\babelfont{tt}[Script=Hebrew]{Miriam Mono CLM} % Hebrew monospace
% Latin-script fonts for embedded English / code / math.
\babelfont[english]{rm}{Latin Modern Roman}
\babelfont[english]{tt}{Latin Modern Mono}
\newfontfamily\codefont{Latin Modern Mono}     % inline English code snippets

% --- Colours -----------------------------------------------------------------
\usepackage{xcolor}
\definecolor{brandblue}{HTML}{1F3B73}
\definecolor{brandred}{HTML}{C0392B}
\definecolor{brandgreen}{HTML}{2E7D57}

% --- Graphics, tables, math --------------------------------------------------
\usepackage{graphicx}
% Figures generated by the Python figure service live in ../assets/figures.
\graphicspath{{../assets/figures/}{figures/}}
\usepackage{booktabs}
\usepackage{array}
\usepackage{amsmath}
\usepackage{amssymb}

% --- TikZ (block/flow diagrams, FR-B8) ---------------------------------------
\usepackage{tikz}
\usetikzlibrary{positioning,arrows.meta,shapes.geometric,calc}

% --- Headers / footers (FR-B2) ----------------------------------------------
\usepackage{fancyhdr}
\pagestyle{fancy}
\fancyhf{}
\fancyhead[R]{\small\textbf{כיצד לבנות סטארט-אפ נכון}}
\fancyhead[L]{\small\nouppercase{\leftmark}}
% Force the page number left-to-right: in the Hebrew (RTL) footer, bare two-digit
% numbers could otherwise be reordered by the bidi algorithm (e.g. "14" -> "41").
\fancyfoot[C]{\small\babelsublr{\thepage}}
\renewcommand{\headrulewidth}{0.4pt}
\renewcommand{\footrulewidth}{0.0pt}
\setlength{\headheight}{15pt}

% Chapter-opening pages use the `plain` style. Redefine it so its page number is
% also forced LTR — the default plain footer reversed two-digit numbers under
% babel RTL (the bug that printed "14" as "41" on a chapter's first page).
\fancypagestyle{plain}{%
  \fancyhf{}%
  \fancyfoot[C]{\small\babelsublr{\thepage}}%
  \renewcommand{\headrulewidth}{0pt}%
  \renewcommand{\footrulewidth}{0pt}%
}

% --- Visual design layer (palette, headings, callout/formula boxes) ----------
% =============================================================================
% theme.tex — visual design layer: extended palette, chapter/section styling,
% callout boxes and highlighted formula boxes. Kept separate from preamble so
% each file stays < 150 LOC. Loaded after xcolor + tikz, before hyperref.
% =============================================================================

\usepackage{titlesec}

% --- BiDi-safe TikZ wrapper --------------------------------------------------
% babel bidi=basic reorders node text ACROSS a multi-node picture (only the
% first node keeps its text). Typesetting the whole picture in an LTR (english)
% context fixes it. Node labels must therefore be English — put the Hebrew
% explanation in the figure caption. Usage: \entikz{ \begin{tikzpicture}... }.
\newcommand{\entikz}[1]{\foreignlanguage{english}{#1}}

% --- Extended palette --------------------------------------------------------
\definecolor{ink}{HTML}{1A2433}        % near-black for sub-headings
\definecolor{panel}{HTML}{EEF2F8}      % soft blue tint (callouts)
\definecolor{panelgreen}{HTML}{E9F3EE} % soft green tint (formulas)
\definecolor{rulegrey}{HTML}{C9D3E0}   % light rule under chapter titles

% --- Chapter banner (titlesec) ----------------------------------------------
% A coloured "פרק N" chip, a large brand-blue title, and a thin accent rule.
% NOTE: under babel RTL the ragged directions are mirrored, so \raggedright
% here flushes the chip + title to the RIGHT margin (the Hebrew reading side).
\titleformat{\chapter}[display]
  {\sffamily\bfseries\raggedright}
  {\colorbox{brandblue}{\color{white}\sffamily\bfseries\large\,פרק~\thechapter\,}}
  {0.7ex}
  {\color{brandblue}\Huge\raggedright}
  [\vspace{0.6ex}{\color{rulegrey}\titlerule[1.2pt]}]
\titlespacing*{\chapter}{0pt}{6pt}{1.4ex}

% --- Section / subsection ----------------------------------------------------
\titleformat*{\section}{\sffamily\bfseries\large\color{brandblue}}
\titleformat*{\subsection}{\sffamily\bfseries\color{ink}}

% --- List markers ------------------------------------------------------------
% David CLM has no U+2022 (•), so the default itemize bullet renders as a blank
% "tofu" box. Use a font-independent, on-brand square rule for the first level.
\renewcommand{\labelitemi}{\textcolor{brandblue}{\raisebox{0.25ex}{\rule{0.6ex}{0.6ex}}}}

% --- Coloured callout / formula boxes ----------------------------------------
% Built from \fcolorbox over a captured savebox (NOT tcolorbox, whose frame
% mis-places under babel RTL). The minipage content inherits RTL direction.
% A savebox captures the body first so minipage's internal grouping can't close
% the frame argument early.
\newsavebox{\tkbox}
\newsavebox{\fmbox}

% \begin{takeaway} ... \end{takeaway}  (optional [custom title]).
\newenvironment{takeaway}[1][שורה תחתונה]
  {\par\smallskip\begin{lrbox}{\tkbox}%
   \begin{minipage}{0.90\linewidth}%
   {\sffamily\bfseries\color{brandblue}#1.}\ }
  {\end{minipage}\end{lrbox}%
   \begin{center}\setlength{\fboxsep}{9pt}%
   \fcolorbox{brandblue}{panel}{\usebox{\tkbox}}\end{center}\smallskip}

% \begin{formulabox} \[ ... \] \end{formulabox}
\newenvironment{formulabox}
  {\par\smallskip\begin{lrbox}{\fmbox}%
   \begin{minipage}{0.60\linewidth}\centering}
  {\end{minipage}\end{lrbox}%
   \begin{center}\setlength{\fboxsep}{8pt}%
   \fcolorbox{brandgreen}{panelgreen}{\usebox{\fmbox}}\end{center}\smallskip}


% --- Bibliography (biblatex + biber, FR-B9) ----------------------------------
\usepackage[autostyle=false,style=english]{csquotes}
\usepackage[backend=biber,style=numeric,sorting=none]{biblatex}
\addbibresource{references.bib}

% --- Hyperlinks (linked citations / TOC) -------------------------------------
\usepackage{hyperref}
\hypersetup{
  colorlinks=true,
  linkcolor=brandblue,
  citecolor=brandgreen,
  urlcolor=brandred,
  pdftitle={How to Correctly Build a Start-up},
  pdfauthor={Yosef Shanaa, Ahmad Kais}
}

% --- Convenience macros ------------------------------------------------------
% \en{...} typesets an embedded English (LTR) span inside Hebrew (RTL) text;
% babel keeps it upright and switches to the Latin font + English hyphenation.
\newcommand{\en}[1]{\foreignlanguage{english}{#1}}
% \term{English term} — emphasised inline English term.
\newcommand{\term}[1]{\foreignlanguage{english}{\textit{#1}}}
